\chapter{绪论}

\section{研究背景与问题来源}

石器时代这类运行时间很长的网络游戏，往往会在玩家群体中沉淀出大量经验说法。它们有的来自长期实战观察，有的来自旧版私服口径，有的则是在不同版本、不同分支、不同翻译文本之间不断转述后形成的“近似共识”。这类经验说法并非没有价值，但它们通常存在三个问题：第一，只描述现象，不明确变量进入公式的顺序；第二，把局部成立的经验外推到全局；第三，忽略取整、上限和状态更新顺序带来的离散跳变。对石器时代而言，这三个问题几乎正好覆盖了宠物养成与战斗系统最容易争议的部分。

本研究直接缘起于一类反复出现的具体问题：宠物转生后的成长上下限为什么会出现经验上“超模”或“不到位”的情况；融合和蛋喂药到底何时取整、何时压缩；战斗中的乱敏、命中、会心、结界、属性反转和骑宠攻防到底看哪一层属性。若只看玩家经验，这些问题往往只能得到零散回答；但若直接回到源码，则有机会把它们统一到同一套变量和状态转移框架下分析。也正因为如此，本文不把战斗系统和宠物养成系统分开研究，而是试图把两者接成一条完整的因果链。

\section{研究目标}

本文的目标可以概括为四点。第一，建立统一的变量与符号体系，使玩家口径、工具口径和源码变量能够一一对应。第二，还原战斗系统中的关键定量机制，包括乱敏、回避、命中、会心、法术命中、场地、结界、属性加强、属性反转以及骑宠伤害。第三，把宠物成长、转生、融合、蛋喂药和孵化过程纳入同一分析框架，并说明它们如何继续影响高等级战斗属性。第四，在机制还原的基础上，把若干典型玩家问题重新表述为可计算的离散优化问题。

更进一步说，本文并不满足于“说明某个公式是什么”，而是希望回答“该公式的变量如何影响结果”“哪些经验判断只在局部成立”“哪些最优化直觉会被取整和上限破坏”。因此，全文会反复强调变量进入顺序、边界条件和分段结构，而不把系统误写成连续线性的近似模型。

\section{研究方法}

本文以当前持有的 \code{gmsv} 源码为主要依据，并以项目 \code{sa_pet_assistant} 中已经对齐源码的工具实现作为辅助验证环境。具体方法包括：

第一，\emph{源码还原}。对关键函数、状态变量和宏展开做逐段阅读，确认公式结构、状态更新顺序和取整位置。第二，\emph{符号化整理}。把散落在不同文件与不同函数中的变量映射到统一符号体系，尽量减少章节之间的口径漂移。第三，\emph{分段函数推导}。对乱敏、命中、会心、伤害和喂药压缩等机制写出可分析的分段函数。第四，\emph{状态转移分析}。对场地、结界、属性加强、属性反转等动态状态，用“写入-读取-衰减-覆盖”的视角理解其演化。第五，\emph{数值与图形辅助}。对部分关键公式生成热力图或切片图，用于展示变量影响与局部非单调现象。第六，\emph{优化建模}。将融合蛋喂药等问题形式化为整数约束下的目标优化问题。

本文采取的总体原则是保守和可验证。对于源码中未明确实现或尚未验证完全的机制，本文只给出保守描述，而不依赖经验去补全“应该如此”的公式。这样做虽然会牺牲部分叙述上的圆满性，但能保证结论与代码口径尽可能一致。

\section{本文结构}

全文结构如下。第 2 章建立研究对象、符号与源码口径；第 3 章和第 4 章分别讨论宠物基础成长模型，以及转生、融合、蛋喂药与孵化链条；第 5 章和第 6 章整理战斗有效属性、状态变量与乱敏机制；第 7 章和第 8 章讨论物理概率与法术命中；第 9 章研究属性系统的动态变换；第 10 章给出骑宠、攻防混合与物理伤害公式；第 11 章把宠物成长映射回战斗能力；第 12 章把相关问题形式化为优化模型并讨论案例；第 13 章总结全文并指出后续工作。

整体上，本文遵循“先统一口径，再还原机制，最后做优化分析”的顺序。这种结构安排的出发点很直接：如果没有前面的变量口径和离散规则，后面的最优化结论就无法判断究竟是在优化源码，还是只是在优化某种近似直觉。

